\chapter{4x10\textsuperscript{9}D.C.}\label{4x10^9}

\hfill\break

Il dolore era difficile da gestire, anche con la morfina che Diane aveva
comprato a un prezzo ridicolo in una farmacia di Padang. La febbre era
anche peggio.

Non era continua. Arrivava a ondate, sciami, bolle di calore e rumore
che mi scoppiavano in testa inaspettatamente. Mi rendeva il corpo
capriccioso, imprevedibile. Una notte cercai a tentoni un bicchiere
d'acqua inesistente e fracassai una lampada da comodino, svegliando la
coppia nella stanza accanto.

L'indomani mattina, recuperata una momentanea lucidità, non ricordavo
l'incidente, ma vidi il sangue rappreso sulle nocche e sentii Diane che
pagava il portiere arrabbiato.

«L'ho fatto davvero?» le chiesi.

«Temo di sì.»

Si sedette su una sedia di vimini accanto al letto. Aveva ordinato il
servizio in camera, uova strapazzate e succo d'arancia, quindi immaginai
che fosse mattina. Oltre le tende trasparenti il cielo era azzurro. La
porta del balcone era aperta e lasciava entrare sbuffi di piacevole aria
calda e l'odore dell'oceano. «Mi dispiace» mi scusai.

«Non eri in te. Vorrei dirti di dimenticartene. Se non fosse che
comunque \emph{l'hai già fatto}.» Mi mise una mano sulla fronte per
tranquillizzarmi. «E non è ancora finita, temo.»

«Da quanto\ldots?»

«È passata una settimana.»

«Solo una settimana?»

«Sì, solo una.»

Non ero nemmeno a metà strada del calvario.

\separatorefiori

Ma gli intervalli di lucidità erano utili per scrivere.

La grafomania era uno dei vari effetti collaterali del farmaco. Diane,
quando si era sottoposta allo stesso tormento, una volta aveva scritto
la frase ``\emph{Sono forse io il custode di mio fratello}'' ripetendola
identica per centinaia di volte su quattordici fogli protocollo.

La mia grafomania almeno era un po' più coerente. Accatastavo pagine
scritte a mano sul comodino mentre attendevo che la febbre lanciasse una
rinnovata offensiva, rileggendo ciò che avevo scritto nel tentativo di
fissarlo nella mente.

Diane trascorse la giornata fuori dell'hotel. Quando tornò le chiesi
dove fosse stata.

«A fare conoscenze.»

Mi disse di aver contattato un intermediario dei trasporti, un uomo
dell'area di Minangkabau di nome Jala la cui attività di import-export
serviva come copertura per quella più redditizia di mediatore per
l'emigrazione.

«Al molo lo conoscono tutti, Jala» mi spiegò. Aveva fatto un'offerta per
dei posti letto in concorrenza con un branco di Kibbutzim, utopisti
scatenati; sebbene l'affare non fosse concluso, si sentiva cautamente
ottimista.

«Fa' attenzione» mi raccomandai. «Potrebbe ancora esserci qualcuno che
ci cerca.»

«Per ora mi pare di no, ma\ldots» Si strinse nelle spalle. Lanciò
un'occhiata al taccuino che avevo in mano. «Scrivi di nuovo?»

«Non mi fa pensare al dolore.»

«Riesci a reggere bene la penna?»

«È come avere l'artrite reumatoide, ma ce la faccio.» ``Per ora''
pensai. «La distrazione vale il fastidio.»

Ma non si trattava solo di quello, ovviamente. Né di semplice
grafomania. Scrivere era un modo per esorcizzare il mio sentirmi in
pericolo.

«Hai fatto un ottimo lavoro. Davvero» disse Diane.

La fissai sconvolto. «Lo hai \emph{letto}?»

«Me lo hai chiesto tu. Mi hai supplicato, Tyler.»

«Deliravo?»

«A quanto pare, anche se in quel momento sembravi piuttosto
ragionevole.»

«Non l'ho scritto pensando che qualcuno lo avrebbe letto.» Ed ero
sconvolto per il fatto che non mi ricordassi di averglielo mostrato.
Quante altre cose potevano già essere svanite?

«Allora non lo guarderò più. Ma quello che hai scritto\ldots» drizzò la
testa. «Sono sorpresa e lusingata che a quel tempo provassi per me
sentimenti tanto forti.»

«Non dovrebbe sorprenderti.»

«Più di quanto credi, invece. Ma è un paradosso, Tyler. La ragazza che
descrivi è indifferente, quasi crudele.»

«Non ho mai pensato che lo fossi.»

«Non è la tua opinione che mi preoccupa. È la mia.»

Mi imponevo di rimanere seduto nel letto, immaginando che fosse un atto
di forza, una prova del mio stoicismo. Più probabilmente, era la prova
che in quel momento gli antidolorifici stavano funzionando. Avevo i
brividi. I brividi erano il primo segnale di un ritorno della febbre.
«Vuoi sapere quando mi sono innamorato di te? Forse dovrei raccontarlo.
È importante. Successe quando avevo dieci anni\ldots»

«Tyler, Tyler. Nessuno s'innamora a dieci anni.»

«È stato quando è morto Sant'Agostino.»

Sant'Agostino era uno springer spaniel con il pedigree, vivace, bianco e
nero, il cucciolo speciale di Diane. Lo aveva ribattezzato San Cane.

Fece una smorfia. «Che cosa macabra.»

Ma io ero serio. E.D. Lawton aveva comprato il cane d'impulso,
probabilmente perché voleva qualcosa con cui ornare il focolare della
Grande Casa, come fosse una coppia di antichi alari. Ma San Cane si era
opposto al suo destino. Era abbastanza ornamentale ma anche curioso e
combinaguai. Con il tempo E.D. arrivò a detestarlo, Carol Lawton lo
ignorava, Jason ne era sconcertato ma in maniera affettuosa. Solo Diane,
che aveva dodici anni, gli si affezionò davvero. Ognuno riusciva a
tirare fuori il meglio dall'altro. Per sei mesi San Cane l'aveva seguita
ovunque tranne che sul minibus della scuola. Nelle sere estive giocavano
insieme sul grande prato, e fu allora che guardai Diane per la prima
volta in modo diverso - la prima volta che provai piacere nel solo
guardarla. Correva con San Cane fino allo sfinimento, e San Cane
aspettava paziente che lei riprendesse fiato. Gli dedicava attenzioni
che nessuno dei Lawton aveva mai neanche provato a dargli - era
sensibile ai suoi stati d'animo e Sant'Agostino la ricambiava.

Non avrei saputo dire perché mi piacesse questo aspetto di lei. Ma nel
mondo inquieto ed emotivamente forte dei Lawton era un'oasi di affetto
senza complicazioni. Se fossi stato un cane, ne sarei stato geloso. E
invece mi colpì che Diane fosse speciale, diversa dalla sua famiglia per
le cose importanti. Affrontava il mondo con un'apertura emozionale che
gli altri Lawton avevano perso o non avevano mai appreso.

Sant'Agostino morì all'improvviso e prematuramente quell'autunno - era
ancora poco più che un cucciolo. Diane era distrutta dal dolore e io
capii di esserne innamorato\ldots{} No, questo sì che sembra macabro.
Non mi innamorai di lei perché piangeva il suo cane. Mi innamorai di lei
perché era \emph{capace} di piangere il suo cane, laddove tutti gli
altri sembravano indifferenti o segretamente sollevati che Sant'Agostino
finalmente se ne fosse andato.

Distolse lo sguardo dal letto volgendolo alla finestra assolata. «La
morte di quella bestiolina mi spezzò il cuore.»

Avevamo seppellito San Cane nel tratto boscoso oltre il prato. Diane
aveva fatto un mucchietto di pietre a mo' di monumento e ogni primavera
lo ricostruiva, fino a quando se ne andò di casa dieci anni dopo.

A ogni cambio di stagione pregava su quei sassi, in silenzio e a mani
giunte. Chi pregasse o per cosa non lo so. Non so cosa faccia la gente
quando prega. Non credo di esserne capace.

Ma per me fu la prima prova che Diane viveva in un mondo ancora più
grande della Grande Casa, un mondo in cui gioia e dolore si spostavano
lentamente come le maree, portandosi dietro il peso di un oceano.

\separatorefiori

Quella notte mi tornò la febbre. Non ricordo niente se non il timore
ricorrente (tornava a intervalli di un'ora) che il farmaco avesse
cancellato più memoria di quanta sarei mai riuscito a recuperare, un
senso di perdita irrimediabile simile a quei sogni in cui si cerca
inutilmente il portafogli scomparso, l'orologio, un oggetto prezioso o
la percezione di se stessi. Mi sembrava di sentire il farmaco marziano
che agiva dentro di me, intento a sferrare nuovi attacchi e a negoziare
tregue momentanee con il mio sistema immunitario, a stabilire teste di
ponte cellulari, a sequestrare sequenze cromosomiche ostili.

Quando tornai in me, Diane non c'era. Protetto contro il dolore dalla
morfina che mi aveva dato, mi alzai dal letto e riuscii a usare il
bagno, poi mi trascinai sul balcone.

Ora di cena. Il Sole era ancora alto ma il cielo era diventato di un
azzurro più scuro. L'aria odorava di latte di cocco e gas di scarico di
motori diesel. L'Arco luccicava a ovest come argento vivo ghiacciato.

Mi ritrovai di nuovo con la voglia di scrivere, il desiderio montava
come un'eco della febbre. Portavo con me il taccuino che avevo riempito
per metà con scarabocchi appena decifrabili. Avrei dovuto chiedere a
Diane di comprarmene un altro. Forse un altro paio. Che poi avrei
riempito di parole.

Parole come ancore, barche di memoria ormeggiate per non colare a picco
con la tempesta.
